% !TEX program = pdflatex
% !TEX encoding = UTF-8 Unicode
% =========================================================
%  TCC -- Trabalho de Conclusão de Curso
%  Norma: ABNT NBR 14724:2011
%  Classe: abnTeX2
% =========================================================

\documentclass[
  12pt,             % tamanho da fonte
  openright,        % capítulos começam em pág ímpar
  twoside,          % impressão em frente e verso
  a4paper,          % tamanho do papel
  english,          % idioma adicional
  brazilian         % idioma principal
]{abntex2}

% -----------------------------------------------------------
% PACOTES
% -----------------------------------------------------------
\usepackage[T1]{fontenc}
\usepackage[utf8]{inputenc}
\usepackage{lmodern}           % fontes Latin Modern
\usepackage{microtype}         % melhor justificação
\usepackage{graphicx}          % inclusão de imagens
\usepackage{float}             % controle de posição de floats
\usepackage{booktabs}          % tabelas profissionais
\usepackage{multirow}          % células multi-linha em tabelas
\usepackage[table]{xcolor}     % cores
\usepackage{listings}          % trechos de código-fonte
\usepackage{caption}           % personalização de legendas
\usepackage{subcaption}        % subfiguras
\usepackage{amsmath}           % matemática
\usepackage{hyperref}          % links e metadados PDF
\usepackage{url}               % URLs
\usepackage{indentfirst}       % indentar 1.º parágrafo de cada seção
\usepackage[alf,abnt-emphasize=bf]{abntex2cite} % citações ABNT

% -----------------------------------------------------------
% CONFIGURAÇÃO DE LISTINGS (CÓDIGO-FONTE)
% -----------------------------------------------------------
\lstset{
  basicstyle=\small\ttfamily,
  breaklines=true,
  frame=single,
  numbers=left,
  numberstyle=\tiny,
  keywordstyle=\bfseries,
  commentstyle=\itshape,
  stringstyle=\ttfamily,
  showstringspaces=false,
  captionpos=b,
  literate=
    {á}{{\'a}}1 {é}{{\'e}}1 {í}{{\'i}}1 {ó}{{\'o}}1 {ú}{{\'u}}1
    {ã}{{\~a}}1 {õ}{{\~o}}1 {â}{{\^a}}1 {ê}{{\^e}}1 {ô}{{\^o}}1
    {à}{{\`a}}1 {ç}{{\c{c}}}1 {Á}{{\'A}}1 {É}{{\'E}}1 {Ã}{{\~A}}1
}

% -----------------------------------------------------------
% CONFIGURAÇÃO DO HYPERREF
% -----------------------------------------------------------
\hypersetup{
  pdftitle={Desenvolvimento de uma Plataforma Web Gamificada para Gestão de Estudos Universitários},
  pdfauthor={[SEU NOME]},
  pdfsubject={Trabalho de Conclusão de Curso},
  pdfkeywords={gamificação, gestão de estudos, Spring Boot, arquitetura hexagonal, REST API},
  colorlinks=true,
  linkcolor=black,
  citecolor=black,
  urlcolor=blue
}

% -----------------------------------------------------------
% INFORMAÇÕES DO TRABALHO (PREENCHA AQUI)
% -----------------------------------------------------------
\titulo{Desenvolvimento de uma Plataforma Web Gamificada para Gestão de Estudos Universitários}
\autor{[SEU NOME COMPLETO]}
\local{[CIDADE]}
\data{2026}
\orientador{[NOME DO ORIENTADOR]}
% \coorientador{[NOME DO COORIENTADOR]}  % descomente se houver
\instituicao{%
  [NOME DA INSTITUIÇÃO]\\
  [FACULDADE / DEPARTAMENTO]
}
\tipotrabalho{Trabalho de Conclusão de Curso}
\preambulo{%
  Trabalho de Conclusão de Curso apresentado ao [FACULDADE/DEPARTAMENTO]
  da [NOME DA INSTITUIÇÃO] como requisito parcial para a obtenção do grau de
  Bacharel em [NOME DO CURSO].%
}

% -----------------------------------------------------------
% INÍCIO DO DOCUMENTO
% -----------------------------------------------------------
\begin{document}

\frenchspacing % espaçamento correto após ponto final

% ==========================================================
% ELEMENTOS PRÉ-TEXTUAIS
% ==========================================================
\pretextual

% Capa (obrigatório)
\imprimircapa

% Folha de rosto (obrigatório)
\imprimirfolhaderosto*  % * = para ficha catalográfica no verso

% -----------------------------------------------------------
% FICHA CATALOGRÁFICA (verso da folha de rosto)
% A ficha é gerada pela biblioteca da sua instituição.
% Substitua pela imagem fornecida pela biblioteca.
% -----------------------------------------------------------
\begin{fichacatalografica}
  \vspace*{\fill}
  \begin{center}
    \fbox{%
      \begin{minipage}[c][8cm]{13cm}
        \small
        \textbf{[SEU NOME COMPLETO]}

        \hspace{0.5cm} \textit{Desenvolvimento de uma Plataforma Web Gamificada para
        Gestão de Estudos Universitários} / [SEU NOME COMPLETO]. --
        [CIDADE], \imprimirdata.

        \hspace{0.5cm} [NNN] p.; 30 cm.

        \hspace{0.5cm} \imprimirorientadorRotulo~\imprimirorientador

        \hspace{0.5cm}
        \parbox[t]{\textwidth}{\imprimirtipotrabalho~--~\imprimirinstituicao,
        \imprimirdata.}

        \hspace{0.5cm}
          1. Gamificação.
          2. Gestão de estudos.
          3. Arquitetura hexagonal.
          4. Spring Boot.
          5. API REST.
          I. \imprimirorientador.
          II. \imprimirinstituicao.
          III. \imprimirtitulo.
      \end{minipage}
    }
  \end{center}
\end{fichacatalografica}

% -----------------------------------------------------------
% FOLHA DE APROVAÇÃO
% Após a defesa, substitua este ambiente pela digitalização
% da folha assinada pela banca.
% -----------------------------------------------------------
\begin{folhadeaprovacao}
  \begin{center}
    {\ABNTEXchapterfont\large\imprimirautor}

    \vspace*{\fill}\vspace*{\fill}

    \begin{center}
      \ABNTEXchapterfont\bfseries\Large\imprimirtitulo
    \end{center}

    \vspace*{\fill}

    \hspace{.45\textwidth}
    \begin{minipage}{.5\textwidth}
      \imprimirpreambulo
    \end{minipage}

    \vspace*{\fill}
  \end{center}

  \center Aprovado em: \underline{\hspace{2cm}} de \underline{\hspace{3cm}} de \underline{\hspace{1.5cm}}.

  \assinatura{\textbf{\imprimirorientador} \\ Orientador}
  \assinatura{\textbf{[NOME DO MEMBRO 2]} \\ [Titulação e Instituição]}
  \assinatura{\textbf{[NOME DO MEMBRO 3]} \\ [Titulação e Instituição]}

  \begin{center}
    \imprimirlocal, \imprimirdata
  \end{center}
\end{folhadeaprovacao}

% Dedicatória (opcional)
\begin{dedicatoria}
  \vspace*{\fill}
  \centering
  \noindent
  \textit{[Escreva aqui sua dedicatória.]}
  \vspace*{\fill}
\end{dedicatoria}

% Agradecimentos (opcional)
\begin{agradecimentos}
  [Escreva aqui seus agradecimentos.]
\end{agradecimentos}

% Epígrafe (opcional)
\begin{epigrafe}
  \vspace*{\fill}
  \begin{flushright}
    \textit{``[Citação inspiradora.]''}\\
    [Autor da citação]
  \end{flushright}
\end{epigrafe}

% -----------------------------------------------------------
% RESUMO EM PORTUGUÊS (obrigatório)
% -----------------------------------------------------------
\setlength{\absparsep}{18pt}
\begin{resumo}
  Este trabalho apresenta o desenvolvimento de uma plataforma \textit{web} gamificada
  para apoio à gestão de estudos de estudantes universitários. A solução, denominada
  \textit{Student App}, oferece funcionalidades de organização acadêmica -- controle de
  disciplinas, períodos, atividades, avaliações e horários -- integradas a um sistema
  de gamificação baseado em pontos de experiência (XP), níveis, moedas virtuais e
  sequências de estudo (\textit{streak}). O \textit{backend} da aplicação foi
  desenvolvido em Java com o \textit{framework} Spring Boot, adotando a Arquitetura
  Hexagonal (\textit{Ports \& Adapters}) para garantir isolamento do domínio de negócio
  em relação à infraestrutura. O sistema expõe uma API REST documentada via Swagger e
  utiliza autenticação stateless com JSON Web Tokens (JWT). Os resultados demonstram
  que a abordagem gamificada tem potencial para aumentar o engajamento dos estudantes
  com a plataforma, promovendo hábitos de estudo mais consistentes.

  \vspace{\onelineskip}
  \noindent
  \textbf{Palavras-chave}: gamificação; gestão de estudos; Spring Boot; arquitetura
  hexagonal; API REST; Java.
\end{resumo}

% -----------------------------------------------------------
% ABSTRACT EM INGLÊS (obrigatório)
% -----------------------------------------------------------
\begin{resumo}[Abstract]
  \begin{otherlanguage*}{english}
    This work presents the development of a gamified web platform for academic
    management support aimed at university students. The solution, named
    \textit{Student App}, provides academic organization features -- subject and period
    management, activities, assessments, and class schedules -- integrated with a
    gamification system based on experience points (XP), levels, virtual coins, and
    study streaks. The application backend was developed in Java using the Spring Boot
    framework, adopting the Hexagonal Architecture (\textit{Ports \& Adapters}) pattern
    to ensure domain isolation from infrastructure concerns. The system exposes a
    REST API documented via Swagger and uses stateless authentication with JSON Web
    Tokens (JWT). Results indicate that the gamified approach has the potential to
    increase student engagement with the platform, promoting more consistent study
    habits.

    \vspace{\onelineskip}
    \noindent
    \textbf{Keywords}: gamification; study management; Spring Boot; hexagonal
    architecture; REST API; Java.
  \end{otherlanguage*}
\end{resumo}

% -----------------------------------------------------------
% LISTAS
% -----------------------------------------------------------
\pdfbookmark[0]{\listfigurename}{lof}
\listoffigures*
\cleardoublepage

\pdfbookmark[0]{\listtablename}{lot}
\listoftables*
\cleardoublepage

% Lista de abreviaturas e siglas
\begin{siglas}
  \item[API]   \textit{Application Programming Interface}
  \item[ABNT]  Associação Brasileira de Normas Técnicas
  \item[AWS]   \textit{Amazon Web Services}
  \item[CORS]  \textit{Cross-Origin Resource Sharing}
  \item[CRUD]  \textit{Create, Read, Update, Delete}
  \item[DTO]   \textit{Data Transfer Object}
  \item[HTTP]  \textit{HyperText Transfer Protocol}
  \item[JPA]   \textit{Java Persistence API}
  \item[JSON]  \textit{JavaScript Object Notation}
  \item[JWT]   \textit{JSON Web Token}
  \item[MVC]   \textit{Model-View-Controller}
  \item[ORM]   \textit{Object-Relational Mapping}
  \item[REST]  \textit{Representational State Transfer}
  \item[S3]    \textit{Simple Storage Service}
  \item[SQL]   \textit{Structured Query Language}
  \item[TCC]   Trabalho de Conclusão de Curso
  \item[UUID]  \textit{Universally Unique Identifier}
  \item[XP]    Pontos de Experiência (\textit{Experience Points})
\end{siglas}

% Sumário
\pdfbookmark[0]{\contentsname}{toc}
\tableofcontents*
\cleardoublepage

% ==========================================================
% ELEMENTOS TEXTUAIS
% ==========================================================
\textual

% =========================================================
%  Capítulo 1 — Introdução
% =========================================================

\chapter{Introdução}
\label{cap:introducao}

O ambiente universitário impõe aos estudantes a necessidade de gerenciar,
simultaneamente, múltiplas disciplinas, avaliações, atividades extracurriculares e
compromissos pessoais. Essa sobrecarga cognitiva frequentemente resulta em baixa
produtividade, procrastinação e abandono dos estudos \cite{[REFERENCIA_1]}.
Ferramentas digitais de organização pessoal têm sido amplamente adotadas como
alternativa para mitigar esses problemas, mas muitas carecem de mecanismos que
promovam o engajamento contínuo do usuário.

A gamificação -- aplicação de elementos e mecânicas de jogos em contextos não
lúdicos -- surge como uma abordagem promissora para aumentar a motivação e a
persistência em plataformas educacionais \cite{[REFERENCIA_2]}. Ao incorporar
sistemas de pontuação, níveis e recompensas ao processo de organização acadêmica,
é possível criar um ciclo de \textit{feedback} positivo que incentiva o estudante
a manter hábitos de estudo regulares.

% -----------------------------------------------------------
\section{Contextualização e Motivação}
\label{sec:contextualizacao}

[Descreva aqui o cenário que motivou a criação do projeto. Inclua dados sobre
a dificuldade dos estudantes universitários em gerir a própria vida acadêmica,
pesquisas sobre procrastinação, evasão escolar, etc.]

% -----------------------------------------------------------
\section{Objetivos}
\label{sec:objetivos}

\subsection{Objetivo Geral}
\label{subsec:objetivo-geral}

Desenvolver uma plataforma \textit{web} gamificada para gestão de estudos
universitários, fornecendo ferramentas de organização acadêmica integradas a
mecânicas de engajamento baseadas em gamificação.

\subsection{Objetivos Específicos}
\label{subsec:objetivos-especificos}

\begin{itemize}
  \item Implementar o módulo de organização acadêmica, contemplando o cadastro e
        gerenciamento de disciplinas, períodos letivos, atividades, avaliações e
        horários de aula;

  \item Projetar e implementar um sistema de gamificação com pontos de experiência
        (XP), níveis progressivos, moedas virtuais e controle de sequência de estudo
        (\textit{streak});

  \item Desenvolver um módulo de sessões de foco (\textit{Focus Sessions}) baseado
        na técnica Pomodoro para recompensar o tempo de estudo dedicado;

  \item Construir o \textit{backend} da aplicação adotando a Arquitetura Hexagonal
        (\textit{Ports \& Adapters}) para garantir o isolamento da lógica de negócio
        da infraestrutura;

  \item Expor os recursos da aplicação por meio de uma API REST segura, documentada
        com Swagger e protegida por autenticação JWT;

  \item Implementar um sistema de notificações para alertar o usuário sobre prazos
        de atividades, limites de faltas e conquistas de gamificação.
\end{itemize}

% -----------------------------------------------------------
\section{Justificativa}
\label{sec:justificativa}

[Explique por que este trabalho é relevante. Relacione a necessidade de ferramentas
de gestão acadêmica com o potencial da gamificação para aumentar o engajamento.
Justifique a escolha das tecnologias (Java/Spring Boot, arquitetura hexagonal).]

% -----------------------------------------------------------
\section{Estrutura do Trabalho}
\label{sec:estrutura}

Este trabalho está organizado em seis capítulos. O Capítulo~\ref{cap:referencial}
apresenta o referencial teórico que fundamenta as escolhas tecnológicas e
metodológicas. O Capítulo~\ref{cap:metodologia} descreve a metodologia de
desenvolvimento adotada e as ferramentas utilizadas. O
Capítulo~\ref{cap:desenvolvimento} detalha a implementação do sistema, cobrindo
arquitetura, modelagem de domínio e principais funcionalidades. O
Capítulo~\ref{cap:resultados} apresenta os resultados obtidos. Por fim, o
Capítulo~\ref{cap:conclusao} conclui o trabalho e aponta direções para trabalhos
futuros.

% =========================================================
%  Capítulo 2 — Referencial Teórico
% =========================================================

\chapter{Referencial Teórico}
\label{cap:referencial}

Este capítulo apresenta os conceitos e tecnologias que fundamentam o desenvolvimento
da plataforma proposta.

% -----------------------------------------------------------
\section{Gamificação}
\label{sec:gamificacao}

O termo \textit{gamificação} (\textit{gamification}) foi popularizado por
\citeonline{[REFERENCIA_DETERDING_2011]} e designa o uso de elementos de \textit{design}
de jogos em contextos que não são jogos, com o objetivo de aumentar o engajamento,
a motivação e o comportamento dos usuários.

\subsection{Elementos de Gamificação}
\label{subsec:elementos-gamificacao}

Os principais elementos gamificados utilizados neste trabalho são:

\begin{itemize}
  \item \textbf{Pontos de Experiência (XP):} representam o progresso do usuário ao
        completar atividades e sessões de estudo. O acúmulo de XP determina o nível
        atual do estudante;

  \item \textbf{Níveis:} estrutura de progressão que requer quantidades crescentes
        de XP. O limiar de XP por nível é calculado como $100 \times \text{nível\_atual}$;

  \item \textbf{Moedas Virtuais:} recompensas adicionais que podem ser utilizadas
        para obter benefícios dentro da plataforma;

  \item \textbf{Sequência de Estudo (\textit{Streak}):} contador de dias consecutivos
        em que o usuário realiza atividades na plataforma, incentivando a consistência.
\end{itemize}

\subsection{Gamificação no Contexto Educacional}
\label{subsec:gamificacao-educacional}

[Apresente pesquisas e estudos sobre o impacto da gamificação em ambientes
educacionais. Cite trabalhos que demonstrem melhoria no engajamento e desempenho
dos estudantes.]

% -----------------------------------------------------------
\section{Gestão de Estudos e Técnica Pomodoro}
\label{sec:gestao-estudos}

[Descreva o problema da gestão de tempo no contexto acadêmico. Apresente a Técnica
Pomodoro -- desenvolvida por Francesco Cirillo -- como método de organização de
sessões de foco com intervalos regulares, e como ela foi integrada ao sistema de
gamificação desta plataforma.]

% -----------------------------------------------------------
\section{Arquitetura de \textit{Software}}
\label{sec:arquitetura-software}

\subsection{Arquitetura Hexagonal (\textit{Ports \& Adapters})}
\label{subsec:arquitetura-hexagonal}

A Arquitetura Hexagonal, proposta por \citeonline{[REFERENCIA_COCKBURN_2005]}, também
conhecida como \textit{Ports \& Adapters}, tem como objetivo central isolar o núcleo
da aplicação (domínio de negócio) de agentes externos, como interfaces de usuário,
bancos de dados e sistemas de mensageria.

O modelo é composto por três camadas principais:

\begin{itemize}
  \item \textbf{Núcleo (Domínio):} contém as regras de negócio, entidades e casos
        de uso, sem qualquer dependência de frameworks ou tecnologias externas;

  \item \textbf{Portas de Entrada (\textit{Input Ports}):} interfaces que definem
        os casos de uso que o domínio expõe para o mundo externo;

  \item \textbf{Portas de Saída (\textit{Output Ports}):} interfaces que definem
        os recursos externos que o domínio necessita (repositórios, serviços de e-mail,
        armazenamento de arquivos etc.).
\end{itemize}

Os \textit{Adapters} implementam essas interfaces, conectando o domínio à tecnologia
concreta. A Figura~\ref{fig:arquitetura-hexagonal} ilustra a estrutura de camadas
adotada neste trabalho.

\begin{figure}[H]
  \centering
  % \includegraphics[width=0.8\textwidth]{figuras/arquitetura-hexagonal.png}
  \fbox{\parbox{0.8\textwidth}{\centering\vspace{2cm}[Inserir diagrama da Arquitetura Hexagonal]\vspace{2cm}}}
  \caption{Arquitetura Hexagonal do Student App}
  \label{fig:arquitetura-hexagonal}
  \fonte{Elaborado pelo autor (\imprimirdata).}
\end{figure}

\subsection{API REST}
\label{subsec:api-rest}

O estilo arquitetural \textit{Representational State Transfer} (REST), descrito por
\citeonline{[REFERENCIA_FIELDING_2000]}, define um conjunto de restrições para o
projeto de sistemas distribuídos baseados em hipermídia. Os princípios REST incluem:
comunicação \textit{stateless}, interface uniforme, uso dos verbos HTTP de forma
semântica e recursos identificados por URIs.

[Continue descrevendo REST e como foi aplicado no projeto.]

% -----------------------------------------------------------
\section{Tecnologias Utilizadas}
\label{sec:tecnologias}

\subsection{Java e Spring Boot}
\label{subsec:spring-boot}

Java é uma linguagem de programação orientada a objetos amplamente adotada no
desenvolvimento de sistemas corporativos \cite{[REFERENCIA_JAVA]}. O framework
Spring Boot, na versão 3.x, simplifica a criação de aplicações Java ao fornecer
convenções de configuração automática, um servidor embutido e um ecossistema rico
de bibliotecas \cite{[REFERENCIA_SPRING]}.

Neste trabalho, utiliza-se o Spring Boot versão 3.5.6 com Java 17 como plataforma
base para o \textit{backend}.

\subsection{Spring Data JPA e PostgreSQL}
\label{subsec:jpa-postgres}

[Descreva a camada de persistência. Explique JPA (Java Persistence API), Hibernate
como implementação ORM, e PostgreSQL como banco de dados relacional escolhido.]

\subsection{MapStruct}
\label{subsec:mapstruct}

[Descreva o MapStruct como gerador de código para mapeamento entre objetos (domain
models → JPA entities e vice-versa), e por que ele foi escolhido em relação a
alternativas como ModelMapper.]

\subsection{JSON Web Tokens (JWT)}
\label{subsec:jwt}

[Descreva JWT como mecanismo de autenticação stateless. Explique a estrutura do
token (header, payload, signature) e como ele é validado a cada requisição no
filtro JwtAuthFilter.]

\subsection{Amazon S3}
\label{subsec:aws-s3}

[Descreva o uso do AWS S3 para armazenamento de arquivos (materiais de estudo)
enviados pelos usuários.]

% -----------------------------------------------------------
\section{Trabalhos Relacionados}
\label{sec:trabalhos-relacionados}

[Apresente e compare aplicações similares como Google Classroom, Notion, Duolingo
(gamificação), Habitica (gamificação para hábitos), etc. Identifique o diferencial
do Student App em relação a essas soluções.]

% =========================================================
%  Capítulo 3 — Metodologia
% =========================================================

\chapter{Metodologia}
\label{cap:metodologia}

Este capítulo descreve a metodologia adotada para o planejamento, desenvolvimento
e validação da plataforma.

% -----------------------------------------------------------
\section{Classificação da Pesquisa}
\label{sec:classificacao-pesquisa}

Do ponto de vista da sua natureza, este trabalho caracteriza-se como uma pesquisa
aplicada, pois tem como objetivo gerar conhecimento para solução de um problema
específico \cite{[REFERENCIA_METODOLOGIA]}. Quanto à abordagem, é predominantemente
qualitativa, uma vez que a avaliação do sistema envolve análise de funcionalidades
e usabilidade. Com relação aos procedimentos técnicos, trata-se de um estudo de
caso de desenvolvimento de \textit{software}.

% -----------------------------------------------------------
\section{Metodologia de Desenvolvimento}
\label{sec:metodologia-desenvolvimento}

O desenvolvimento foi conduzido de forma iterativa e incremental, com ênfase
em entregas funcionais a cada ciclo. As funcionalidades foram organizadas em
etapas:

\begin{enumerate}
  \item \textbf{Modelagem do domínio:} levantamento e especificação das entidades,
        regras de negócio e casos de uso;

  \item \textbf{Implementação do \textit{backend}:} desenvolvimento em Java/Spring Boot
        seguindo a Arquitetura Hexagonal, com cobertura de testes unitários;

  \item \textbf{Integração e testes:} validação dos componentes por meio de testes
        unitários da camada de domínio (modelos e serviços) e testes de controladores;

  \item \textbf{Documentação:} geração automática da documentação da API via Swagger/
        OpenAPI e elaboração deste documento.
\end{enumerate}

% -----------------------------------------------------------
\section{Ferramentas e Ambiente de Desenvolvimento}
\label{sec:ferramentas}

O Quadro~\ref{qdr:ferramentas} lista as principais ferramentas e tecnologias
utilizadas no desenvolvimento da plataforma.

\begin{table}[H]
  \centering
  \caption{Ferramentas e tecnologias utilizadas}
  \label{qdr:ferramentas}
  \begin{tabular}{lll}
    \toprule
    \textbf{Categoria}        & \textbf{Tecnologia/Ferramenta}    & \textbf{Versão} \\
    \midrule
    Linguagem                 & Java                              & 17              \\
    Framework                 & Spring Boot                       & 3.5.6           \\
    Persistência              & Spring Data JPA / Hibernate       & 6.x             \\
    Banco de dados            & PostgreSQL                        & 15              \\
    Mapeamento objeto-relacional & MapStruct                      & 1.5.x           \\
    Autenticação              & JWT (JJWT)                        & 0.12.x          \\
    Armazenamento             & AWS S3 (SDK v2)                   & 2.x             \\
    Conteinerização           & Docker / Docker Compose           & 24.x            \\
    Build                     & Apache Maven                      & 3.x (wrapper)   \\
    Documentação de API       & Springdoc OpenAPI (Swagger UI)    & 2.x             \\
    IDE                       & IntelliJ IDEA                     & --              \\
    Controle de versão        & Git / GitHub                      & --              \\
    Testes unitários          & JUnit 5 + Mockito                 & --              \\
    \bottomrule
  \end{tabular}
  \fonte{Elaborado pelo autor (\imprimirdata).}
\end{table}

% -----------------------------------------------------------
\section{Estratégia de Testes}
\label{sec:estrategia-testes}

A estratégia de testes adotada concentra-se na camada de domínio, em conformidade
com os princípios da Arquitetura Hexagonal, que recomenda que o núcleo da aplicação
seja testável de forma isolada, sem dependência de frameworks ou infraestrutura.

Os testes foram implementados com JUnit 5 e Mockito, cobrindo:

\begin{itemize}
  \item \textbf{Testes de modelos de domínio:} verificação das regras de negócio
        encapsuladas nos métodos das entidades, como \texttt{User.awardXp()} e
        os construtores via \textit{factory methods};

  \item \textbf{Testes de serviços (casos de uso):} verificação do comportamento
        dos \texttt{ServiceImpl} com repositórios \textit{mockados}, garantindo
        o fluxo correto de orquestração entre domínio e portas de saída;

  \item \textbf{Testes de controladores:} verificação dos \textit{endpoints} REST
        com \texttt{@WebMvcTest}, cobrindo casos de sucesso, validações de entrada
        e tratamento de erros HTTP.
\end{itemize}

Testes de integração (\texttt{@SpringBootTest}, \texttt{@DataJpaTest}) estão fora
do escopo deste trabalho.

% -----------------------------------------------------------
\section{Controle de Versão}
\label{sec:controle-versao}

O código-fonte foi versionado com Git, hospedado em repositório privado no GitHub.
O fluxo de trabalho adotado foi baseado em \textit{commits} incrementais organizados
por funcionalidade, com mensagens no formato convencional
(\textit{feat}, \textit{fix}, \textit{refactor}).

% =========================================================
%  Capítulo 4 — Desenvolvimento
% =========================================================

\chapter{Desenvolvimento}
\label{cap:desenvolvimento}

Este capítulo descreve as decisões de projeto e a implementação do \textit{backend}
da plataforma \textit{Student App}.

% -----------------------------------------------------------
\section{Visão Geral do Sistema}
\label{sec:visao-geral}

O \textit{Student App} é uma plataforma \textit{web} voltada para estudantes
universitários que desejam organizar sua vida acadêmica de forma gamificada. O
sistema é composto por um \textit{backend} em Java/Spring Boot que expõe uma API
REST consumida por um \textit{frontend} (fora do escopo deste trabalho).

A Figura~\ref{fig:visao-geral} apresenta a arquitetura de alto nível da solução.

\begin{figure}[H]
  \centering
  \fbox{\parbox{0.85\textwidth}{\centering\vspace{2cm}[Inserir diagrama de visão geral do sistema]\vspace{2cm}}}
  \caption{Visão geral da arquitetura do Student App}
  \label{fig:visao-geral}
  \fonte{Elaborado pelo autor (\imprimirdata).}
\end{figure}

% -----------------------------------------------------------
\section{Arquitetura do Sistema}
\label{sec:arquitetura}

\subsection{Estrutura de Pacotes}
\label{subsec:estrutura-pacotes}

A estrutura de pacotes do projeto reflete diretamente a separação de camadas
da Arquitetura Hexagonal:

\begin{lstlisting}[language=,caption={Estrutura de pacotes do projeto},label={lst:estrutura-pacotes}]
com.studentapp.api
  domain/
    model/          -- Entidades de domínio puras (sem anotações JPA)
    port/in/        -- Interfaces de casos de uso (portas de entrada)
    port/out/       -- Interfaces de repositórios e serviços (portas de saída)
    service/        -- Implementações dos casos de uso (XxxServiceImpl)
  infra/
    adapters/in/web/          -- Controllers REST, DTOs e mapeadores DTO
    adapters/out/persistance/ -- Entidades JPA, repositórios JPA,
                              --   mapeadores de persistência, adaptadores
    adapters/out/email/       -- Adaptador JavaMail
    adapters/out/storage/     -- Adaptador AWS S3
    config/                   -- Segurança, JWT, S3, CORS, tratamento de erros
\end{lstlisting}

\subsection{Fluxo de Requisição}
\label{subsec:fluxo-requisicao}

O fluxo de processamento de uma requisição segue a sequência ilustrada na
Figura~\ref{fig:fluxo-requisicao}:

\[
  \text{Controller} \rightarrow \text{UseCase (porta/in)} \rightarrow
  \text{ServiceImpl} \rightarrow \text{RepositoryPort (porta/out)}
  \rightarrow \text{RepositoryAdapter} \rightarrow \text{JpaRepository}
\]

\begin{figure}[H]
  \centering
  \fbox{\parbox{0.9\textwidth}{\centering\vspace{2cm}[Inserir diagrama de sequência do fluxo de requisição]\vspace{2cm}}}
  \caption{Fluxo de processamento de uma requisição REST}
  \label{fig:fluxo-requisicao}
  \fonte{Elaborado pelo autor (\imprimirdata).}
\end{figure}

% -----------------------------------------------------------
\section{Modelagem do Domínio}
\label{sec:modelagem-dominio}

\subsection{Entidades Principais}
\label{subsec:entidades}

O domínio da aplicação é composto por dezoito entidades, conforme o Diagrama de
Classes do Apêndice~\ref{apendice:diagrama-classes}. O Quadro~\ref{qdr:entidades}
apresenta as entidades e suas responsabilidades.

\begin{table}[H]
  \centering
  \caption{Entidades do domínio e suas responsabilidades}
  \label{qdr:entidades}
  \begin{tabular}{ll}
    \toprule
    \textbf{Entidade}       & \textbf{Responsabilidade}                                   \\
    \midrule
    \texttt{User}           & Usuário do sistema; gerencia gamificação (XP, nível, moedas) \\
    \texttt{UserPreference} & Preferências e configurações do usuário                      \\
    \texttt{Period}         & Período letivo (semestre/ano)                                \\
    \texttt{Subject}        & Disciplina cursada em um período                             \\
    \texttt{ClassSchedule}  & Horário de aula de uma disciplina                            \\
    \texttt{Activity}       & Atividade com \textit{checklist} e prazo                     \\
    \texttt{ChecklistItem}  & Item individual da lista de tarefas de uma atividade         \\
    \texttt{Assessment}     & Avaliação/prova de uma disciplina                            \\
    \texttt{AbsenceLog}     & Registro de falta em uma disciplina                          \\
    \texttt{Note}           & Nota de estudo em texto livre                                \\
    \texttt{Material}       & Material de estudo (arquivo ou \textit{link})                \\
    \texttt{FileObject}     & Metadados de arquivo armazenado no AWS S3                    \\
    \texttt{FocusSession}   & Sessão de foco com duração e XP ganho                        \\
    \texttt{PlannerEvent}   & Evento no planejador do estudante                            \\
    \texttt{Reminder}       & Lembrete pessoal                                             \\
    \texttt{Notification}   & Notificação do sistema para o usuário                        \\
    \texttt{PasswordResetToken} & Token para redefinição de senha                         \\
    \texttt{GamificationConfig} & Configurações globais de gamificação                   \\
    \bottomrule
  \end{tabular}
  \fonte{Elaborado pelo autor (\imprimirdata).}
\end{table}

\subsection{Padrão de Construção das Entidades}
\label{subsec:padrao-construcao}

As entidades de domínio utilizam o padrão \textit{factory method} estático com
construtores privados, garantindo a integridade dos objetos e evitando estados
inválidos. Dois \textit{factory methods} são definidos para cada entidade:

\begin{lstlisting}[language=Java, caption={Factory methods da entidade User},label={lst:factory-methods}]
public class User {
    // Construtor privado: impede instanciação direta
    private User() {}

    // Criação de nova instância: gera UUID e timestamps
    public static User create(String name, String email,
                              String passwordHash, Role role) { ... }

    // Reconstituição a partir da persistência
    public static User fromState(UUID id, String name, String email,
                                 String passwordHash, ...) { ... }
}
\end{lstlisting}

% -----------------------------------------------------------
\section{Sistema de Gamificação}
\label{sec:gamificacao-impl}

\subsection{Progressão de XP e Níveis}
\label{subsec:progressao-xp}

O sistema de gamificação utiliza uma progressão de nível baseada em um limiar
crescente de XP. A equação~\ref{eq:xp-threshold} define a quantidade de XP
necessária para avançar do nível atual para o próximo:

\begin{equation}
  \label{eq:xp-threshold}
  \text{XP}_{\text{necessário}} = 100 \times \text{nível\_atual}
\end{equation}

O método \texttt{User.awardXp(int amount)}, apresentado no
Código~\ref{lst:award-xp}, implementa a lógica de concessão de XP com preservação
de XP excedente (\textit{overflow}) ao subir de nível:

\begin{lstlisting}[language=Java, caption={Método awardXp da entidade User},label={lst:award-xp}]
public boolean awardXp(int amount) {
    this.currentXp += amount;
    int threshold = 100 * this.currentLevel;
    if (this.currentXp >= threshold) {
        this.currentXp -= threshold;  // preserva overflow
        this.currentLevel++;
        touch();
        return true;  // indica subida de nível
    }
    touch();
    return false;
}
\end{lstlisting}

\subsection{Eventos de Gamificação}
\label{subsec:eventos-gamificacao}

A concessão de XP ocorre em três eventos principais:

\begin{itemize}
  \item \textbf{Conclusão de Atividade:} ao marcar uma \texttt{Activity} como
        concluída pela primeira vez, o usuário recebe 50 XP. O
        \texttt{ActivityServiceImpl} verifica se a atividade não estava
        previamente concluída para evitar concessão duplicada;

  \item \textbf{Sessão de Foco:} ao encerrar uma \texttt{FocusSession},
        o XP é proporcional à duração da sessão (\texttt{durationSeconds}),
        calculado via \texttt{GamificationConfig};

  \item \textbf{Notificação de Subida de Nível:} quando \texttt{awardXp()}
        retorna \texttt{true}, uma \texttt{Notification} do tipo
        \texttt{LEVEL\_UP} é criada automaticamente.
\end{itemize}

\subsection{Controle de Faltas e Notificações}
\label{subsec:controle-faltas}

O \texttt{AbsenceLogServiceImpl} monitora o percentual de faltas em relação ao
máximo permitido pela disciplina. Quando o limite de 75\% é atingido, uma
\texttt{Notification} do tipo \texttt{ABSENCE\_LIMIT} é criada, com deduplicação
via \texttt{NotificationRepositoryPort.existsUnreadByUserIdAndTypeAndReferenceId()}.

% -----------------------------------------------------------
\section{Segurança}
\label{sec:seguranca}

\subsection{Autenticação JWT}
\label{subsec:auth-jwt}

A autenticação é implementada com JWT (\textit{JSON Web Token}) em modelo
\textit{stateless}. O fluxo de autenticação é:

\begin{enumerate}
  \item O cliente realiza \textit{login} em \texttt{POST /auth/login} com
        \textit{e-mail} e senha;
  \item O \texttt{JwtService} gera um token assinado com a chave secreta
        configurada em \texttt{jwt.secret};
  \item Nas requisições subsequentes, o cliente inclui o token no cabeçalho
        \texttt{Authorization: Bearer <token>};
  \item O filtro \texttt{JwtAuthFilter} valida o token e carrega o
        \texttt{UserDetails} para cada requisição.
\end{enumerate}

\subsection{Autorização}
\label{subsec:autorizacao}

O controle de acesso utiliza \textit{roles} definidas no enum \texttt{Role}:

\begin{itemize}
  \item \textbf{USER:} acesso aos próprios recursos acadêmicos e de gamificação;
  \item \textbf{ADMIN:} acesso às rotas \texttt{/admin/**} para operações
        administrativas.
\end{itemize}

% -----------------------------------------------------------
\section{Módulos Funcionais}
\label{sec:modulos}

[Descreva cada módulo funcional com mais detalhe: Organização Acadêmica
(Subject, Period, ClassSchedule, Assessment, AbsenceLog), Atividades com Checklist,
Sessões de Foco (Pomodoro), Materiais e Notas, Planejador, Notificações.
Inclua trechos de código relevantes e diagramas de sequência para os fluxos
mais importantes.]

% -----------------------------------------------------------
\section{Camada de Persistência}
\label{sec:persistencia}

\subsection{Entidades JPA e Mapeamento}
\label{subsec:entidades-jpa}

A camada de persistência é implementada com Spring Data JPA (Hibernate 6) e
PostgreSQL. As entidades JPA espelham as entidades de domínio, porém com anotações
de mapeamento ORM. Os mapeadores de persistência, gerados pelo MapStruct, convertem
entre os dois modelos.

Um caso especial é o mapeamento da lista de \texttt{ChecklistItem} na entidade
\texttt{Activity}: a lista é serializada diretamente como JSON na coluna
\texttt{checklist} usando a anotação do Hibernate 6:

\begin{lstlisting}[language=Java, caption={Mapeamento JSON de ChecklistItem},label={lst:json-mapping}]
@JdbcTypeCode(SqlTypes.JSON)
@Column(columnDefinition = "jsonb")
private List<ChecklistItem> checklist;
\end{lstlisting}

\subsection{Diagrama Entidade-Relacionamento}
\label{subsec:der}

O Diagrama Entidade-Relacionamento completo do banco de dados PostgreSQL está
disponível no Apêndice~\ref{apendice:der}.

% -----------------------------------------------------------
\section{API REST e Documentação}
\label{sec:api-rest}

A API expõe [NÚMERO] \textit{endpoints} organizados em [NÚMERO] grupos de recursos.
O Quadro~\ref{qdr:endpoints-principais} lista os principais grupos de recursos.

\begin{table}[H]
  \centering
  \caption{Grupos de recursos da API REST}
  \label{qdr:endpoints-principais}
  \begin{tabular}{ll}
    \toprule
    \textbf{Prefixo}          & \textbf{Recurso}            \\
    \midrule
    \texttt{/auth}            & Autenticação e registro      \\
    \texttt{/users}           & Gerenciamento de usuários    \\
    \texttt{/periods}         & Períodos letivos             \\
    \texttt{/subjects}        & Disciplinas                  \\
    \texttt{/class-schedules} & Horários de aula             \\
    \texttt{/activities}      & Atividades com \textit{checklist} \\
    \texttt{/assessments}     & Avaliações                   \\
    \texttt{/absence-logs}    & Registros de falta           \\
    \texttt{/focus-sessions}  & Sessões de foco              \\
    \texttt{/notes}           & Notas de estudo              \\
    \texttt{/materials}       & Materiais de estudo          \\
    \texttt{/planner-events}  & Eventos do planejador        \\
    \texttt{/notifications}   & Notificações                 \\
    \bottomrule
  \end{tabular}
  \fonte{Elaborado pelo autor (\imprimirdata).}
\end{table}

A documentação interativa da API está disponível via Swagger UI no caminho
\texttt{/swagger-ui.html}, sem necessidade de autenticação.

% =========================================================
%  Capítulo 5 — Resultados
% =========================================================

\chapter{Resultados}
\label{cap:resultados}

Este capítulo apresenta os resultados obtidos com o desenvolvimento da plataforma,
abrangendo os testes realizados, as funcionalidades entregues e uma avaliação
qualitativa do sistema.

% -----------------------------------------------------------
\section{Funcionalidades Implementadas}
\label{sec:funcionalidades}

Ao final do desenvolvimento, as seguintes funcionalidades foram entregues:

\begin{itemize}
  \item \textbf{Autenticação e autorização:} cadastro de usuários, \textit{login}
        com JWT, redefinição de senha por e-mail e controle de acesso baseado em
        \textit{roles};

  \item \textbf{Organização acadêmica:} gerenciamento completo de períodos letivos,
        disciplinas, horários de aula, avaliações e registros de falta;

  \item \textbf{Atividades com \textit{checklist}:} criação e conclusão de
        atividades com lista de tarefas embutida, serializada como JSON;

  \item \textbf{Sessões de foco:} registro de sessões de estudo baseadas no método
        Pomodoro, com concessão de XP proporcional ao tempo dedicado;

  \item \textbf{Notas e materiais:} criação de notas de estudo e upload de
        materiais (arquivos) com armazenamento no AWS S3;

  \item \textbf{Planejador e lembretes:} criação de eventos e lembretes pessoais;

  \item \textbf{Gamificação:} sistema de XP, níveis, moedas e \textit{streak} com
        progressão baseada em ações do usuário;

  \item \textbf{Notificações:} alertas automáticos para atividades próximas do
        prazo, limite de faltas e subida de nível, com deduplicação.
\end{itemize}

% -----------------------------------------------------------
\section{Testes Realizados}
\label{sec:testes}

\subsection{Cobertura dos Testes Unitários}
\label{subsec:cobertura-testes}

[Apresente aqui os resultados dos testes unitários. Inclua uma tabela com as
classes testadas, número de casos de teste e a cobertura obtida. Execute
\texttt{./mvnw test} e insira o relatório.]

\begin{table}[H]
  \centering
  \caption{Resumo dos testes unitários implementados}
  \label{qdr:testes}
  \begin{tabular}{llc}
    \toprule
    \textbf{Classe testada}          & \textbf{Tipo}      & \textbf{N.º de casos} \\
    \midrule
    \texttt{UserTest}                & Modelo de domínio  & [N]                   \\
    \texttt{UserServiceImplTest}     & Serviço            & [N]                   \\
    \texttt{NotificationServiceImplTest} & Serviço        & [N]                   \\
    \texttt{UserControllerTest}      & Controlador        & [N]                   \\
    [Demais classes\ldots]           & --                 & --                    \\
    \bottomrule
  \end{tabular}
  \fonte{Elaborado pelo autor (\imprimirdata).}
\end{table}

\subsection{Casos de Teste Relevantes}
\label{subsec:casos-teste}

[Descreva os casos de teste mais importantes, especialmente os que validam
regras de negócio críticas como o sistema de XP/nível e a deduplicação de
notificações.]

% -----------------------------------------------------------
\section{Avaliação da Arquitetura}
\label{sec:avaliacao-arquitetura}

A adoção da Arquitetura Hexagonal trouxe benefícios observáveis durante o
desenvolvimento:

\begin{itemize}
  \item \textbf{Testabilidade:} os serviços de domínio podem ser testados de
        forma completamente isolada, sem necessidade de instanciar \textit{beans}
        do Spring, banco de dados ou outros componentes de infraestrutura;

  \item \textbf{Substituição de adapters:} a implementação de armazenamento
        pode ser trocada de AWS S3 para outro provedor sem alteração da lógica
        de domínio, bastando implementar a interface \texttt{FileStorageServicePort};

  \item \textbf{Legibilidade:} a separação clara de responsabilidades facilita
        a localização e manutenção do código, especialmente para novos
        contribuidores.
\end{itemize}

Como contrapartida, a arquitetura exige a criação de um número maior de classes
(duas camadas de mapeamento, interfaces de porta, adaptadores), o que aumenta
o \textit{boilerplate} inicial.

% -----------------------------------------------------------
\section{Discussão}
\label{sec:discussao}

[Discuta os resultados obtidos em relação aos objetivos definidos no
Capítulo~\ref{cap:introducao}. Analise pontos fortes e limitações da
solução. Relacione com o referencial teórico apresentado no
Capítulo~\ref{cap:referencial}.]

% =========================================================
%  Capítulo 6 — Conclusão
% =========================================================

\chapter{Conclusão}
\label{cap:conclusao}

Este trabalho apresentou o desenvolvimento do \textit{Student App}, uma plataforma
\textit{web} gamificada para apoio à gestão de estudos universitários. A solução
integra ferramentas de organização acadêmica -- gerenciamento de disciplinas,
períodos, atividades, avaliações e horários -- com mecânicas de engajamento
baseadas em gamificação, como pontos de experiência, níveis progressivos, moedas
virtuais e sequências de estudo.

O \textit{backend} foi desenvolvido em Java 17 com o \textit{framework} Spring Boot
3.5.6, adotando a Arquitetura Hexagonal como padrão estrutural. Essa escolha
arquitetural demonstrou-se adequada para o contexto do projeto, pois garantiu o
isolamento das regras de negócio da infraestrutura, facilitando os testes unitários
e tornando os componentes de infraestrutura (banco de dados, armazenamento de
arquivos, e-mail) substituíveis de forma independente.

Os objetivos específicos definidos na Seção~\ref{subsec:objetivos-especificos}
foram alcançados:

\begin{itemize}
  \item O módulo de organização acadêmica foi implementado, contemplando todas as
        entidades e seus fluxos de CRUD;

  \item O sistema de gamificação com XP, níveis, moedas e \textit{streak} foi
        integrado às ações do usuário;

  \item As sessões de foco com recompensa de XP proporcional ao tempo dedicado
        foram implementadas;

  \item A Arquitetura Hexagonal foi aplicada de forma consistente em todo o projeto;

  \item A API REST foi documentada via Swagger UI e protegida com JWT;

  \item O sistema de notificações com deduplicação para prazos, limites de faltas
        e conquistas foi implementado.
\end{itemize}

% -----------------------------------------------------------
\section{Limitações}
\label{sec:limitacoes}

[Descreva as limitações do trabalho atual, por exemplo: ausência do frontend,
ausência de testes de integração, falta de avaliação com usuários reais, etc.]

% -----------------------------------------------------------
\section{Trabalhos Futuros}
\label{sec:trabalhos-futuros}

Como desdobramentos naturais deste trabalho, destacam-se as seguintes direções
para trabalhos futuros:

\begin{itemize}
  \item \textbf{Desenvolvimento do \textit{frontend}:} implementação da interface
        do usuário em uma tecnologia moderna como React ou Vue.js, consumindo a
        API REST existente;

  \item \textbf{Avaliação com usuários reais:} realização de um estudo de usabilidade
        e de eficácia da gamificação com estudantes universitários, coletando métricas
        de engajamento e satisfação;

  \item \textbf{Sistema de conquistas (\textit{badges}):} expansão do sistema de
        gamificação com conquistas desbloqueáveis por marcos específicos (e.g., primeira
        semana consecutiva de estudo, 100 sessões de foco concluídas);

  \item \textbf{Ranking e elementos sociais:} introdução de tabelas de classificação
        e funcionalidades colaborativas entre estudantes;

  \item \textbf{Análise de desempenho:} implementação de relatórios e visualizações
        sobre o histórico de estudos, distribuição de tempo por disciplina e evolução
        das notas;

  \item \textbf{Aplicativo móvel:} desenvolvimento de clientes móveis (iOS/Android)
        com suporte a notificações \textit{push};

  \item \textbf{Testes de integração e de carga:} ampliação da suíte de testes com
        testes de integração (\texttt{@DataJpaTest}) e testes de desempenho sob carga.
\end{itemize}


% ==========================================================
% ELEMENTOS PÓS-TEXTUAIS
% ==========================================================
\postextual

% Referências (obrigatório)
\bibliography{referencias}

% Glossário (opcional)
% \glossary

% Apêndices (opcional — material do próprio autor)
\begin{apendicesenv}
  \partapendices
  \chapter{Diagrama de Classes do Domínio}
  \label{apendice:diagrama-classes}

  [Insira aqui o diagrama de classes UML do domínio da aplicação.]

  \chapter{Endpoints da API REST}
  \label{apendice:endpoints}

  [Insira aqui a documentação completa dos endpoints. A documentação interativa
  está disponível em \texttt{/swagger-ui.html} quando a aplicação está em execução.]

  \chapter{Diagrama Entidade-Relacionamento}
  \label{apendice:der}

  [Insira aqui o DER do banco de dados PostgreSQL.]

\end{apendicesenv}

% Anexos (opcional — material de terceiros)
% \begin{anexosenv}
%   \partanexos
%   \chapter{...}
% \end{anexosenv}

\end{document}
